%!TEX root = ../main.tex
\chapter{Implementacja}
\label{ch:implementacja}

\section{Środowisko i narzędzia}
\label{sec:srodowisko}

Prototyp zaimplementowano w języku Go (minimalna wersja języka zadeklarowana w pliku \texttt{go.mod} to 1.25), wybranym z przyczyn omówionych w sekcji~\ref{sec:wybor-technologii}. Wykorzystano następujące podstawowe komponenty i narzędzia:

\begin{itemize}
    \item \textbf{CometBFT} w wersji 0.38 (ABCI 2.0) -- silnik konsensusu, dołączony jako zależność biblioteczna; aplikacja osadza go w procesie węzła i komunikuje się z nim lokalnie przez interfejs ABCI;
    \item \textbf{Protocol Buffers} -- schemat danych zdefiniowano w plikach \texttt{.proto}, a kod struktur Go wygenerowano kompilatorem \texttt{protoc} z wtyczką \texttt{protoc-gen-go}; generacja jest zautomatyzowana celem \texttt{make proto-gen};
    \item \textbf{Cobra} -- biblioteka do budowy interfejsu wiersza poleceń klienta;
    \item \textbf{golangci-lint} oraz \texttt{gofumpt} -- statyczna analiza i formatowanie kodu, uruchamiane celami \texttt{make lint} i \texttt{make fmt};
    \item wbudowany w język \textbf{framework testowy} \texttt{go test}, uzupełniony o testy z wykrywaniem wyścigów (\emph{race detector}) oraz testy rozmyte (\emph{fuzz testing}).
\end{itemize}

Budowanie i typowe operacje deweloperskie ujęto w pliku \texttt{Makefile}: cel \texttt{make build} kompiluje trzy programy wykonywalne -- węzeł \texttt{dpkid}, klienta \texttt{dpki-cli} oraz narzędzie pomiarowe \texttt{benchmark} -- do katalogu \texttt{bin/}, a cele \texttt{make test} i \texttt{make test-cover} uruchamiają zestaw testów odpowiednio z wykrywaniem wyścigów i z pomiarem pokrycia.

\section{Struktura projektu}
\label{sec:struktura-projektu}

Kod zorganizowano zgodnie z konwencjami ekosystemu Go, oddzielając programy wykonywalne (katalog \texttt{cmd/}) od pakietów wewnętrznych zawierających logikę (katalog \texttt{internal/}). Podział ten odzwierciedla trójwarstwową architekturę z sekcji~\ref{sec:architektura}.

\begin{itemize}
    \item \texttt{cmd/dpkid} -- punkt wejścia węzła; podkomendy \texttt{init} (inicjalizacja pojedynczego węzła), \texttt{start} (uruchomienie) oraz \texttt{testnet} (generowanie sieci wielowęzłowej, sekcja~\ref{sec:siec-wielowezlowa});
    \item \texttt{cmd/dpki-cli} -- klient; podkatalog \texttt{commands} zawiera definicje komend Cobry, a \texttt{client} -- nakładkę na interfejs RPC węzła;
    \item \texttt{cmd/benchmark} -- narzędzie generujące obciążenie i mierzące wydajność, wykorzystane w rozdziale~\ref{ch:eksperymenty};
    \item \texttt{internal/types} -- model danych: struktury wygenerowane z Protocol Buffers (pliki \texttt{.pb.go}) oraz logika towarzysząca -- konstruktory, walidacja składniowa, wyznaczanie bajtów do podpisania i podpisywanie;
    \item \texttt{internal/crypto} -- prymitywy kryptograficzne: generowanie i obsługa kluczy Ed25519, podpis i weryfikacja oraz funkcja skrótu SHA-256;
    \item \texttt{internal/state} -- warstwa stanu: rzadkie drzewo Merkle, abstrakcja magazynu klucz--wartość, kodowanie dowodów oraz odwzorowanie domeny na wpis stanu;
    \item \texttt{internal/app} -- aplikacja ABCI spinająca powyższe warstwy w deterministyczną maszynę stanów;
    \item \texttt{internal/verifier} -- weryfikacja własności domeny metodą DNS-01;
    \item \texttt{internal/testhelpers} -- współdzielone funkcje pomocnicze testów.
\end{itemize}

\section{Aplikacja ABCI}
\label{sec:aplikacja-abci}

Sercem systemu jest aplikacja ABCI w pakiecie \texttt{internal/app}, reprezentowana strukturą \texttt{App}. Osadza ona typ \texttt{abci.BaseApplication} z biblioteki CometBFT, dzięki czemu metody interfejsu ABCI nieistotne dla DPKI (obsługa migawek stanu, rozszerzeń głosów) zachowują domyślne, puste zachowanie, a nadpisane zostają jedynie metody niosące logikę rejestru: \texttt{Info}, \texttt{InitChain}, \texttt{CheckTx}, \texttt{ProcessProposal}, \texttt{FinalizeBlock}, \texttt{Commit} oraz \texttt{Query}. Strukturę \texttt{App} przedstawia listing~\ref{lst:app-struct}.

\begin{lstlisting}[language=Go,caption={Struktura aplikacji ABCI},label={lst:app-struct}]
type App struct {
    abci.BaseApplication

    mu       sync.RWMutex
    smt      *state.SMT       // stan: rzadkie drzewo Merkle
    verifier verifier.Verifier // weryfikacja wlasnosci domeny
    chainID  string           // identyfikator lancucha z genezy
    height   int64
    logger   *slog.Logger
}
\end{lstlisting}

Dwie decyzje widoczne w tej strukturze wymagają komentarza. Po pierwsze, identyfikator łańcucha (\texttt{chainID}) jest przekazywany aplikacji przy jej tworzeniu, na podstawie pliku genezy, a nie odczytywany z wywołania \texttt{InitChain}. Wywołanie to następuje bowiem tylko raz, przy pierwszym uruchomieniu łańcucha; gdyby aplikacja zależała od niego, po restarcie węzła i ponownym odtwarzaniu bloków przez \texttt{FinalizeBlock} utraciłaby informację potrzebną do walidacji. Po drugie, dostęp do drzewa stanu jest serializowany muteksem \texttt{mu}, ponieważ CometBFT obsługuje połączenia konsensusu, mempoola i zapytań współbieżnie: \texttt{FinalizeBlock} i \texttt{Commit} pobierają blokadę zapisu, a \texttt{CheckTx}, \texttt{Query} oraz \texttt{Info} -- blokadę odczytu.

\subsection{Obsługa CheckTx i FinalizeBlock}
\label{subsec:checktx-finalize}

Najistotniejszym elementem implementacji jest realizacja dwuetapowej walidacji zapowiedzianej w sekcji~\ref{sec:protokoly}. Granica determinizmu wyznaczona przez interfejs ABCI (sekcja~\ref{sec:abci}) przekłada się bezpośrednio na podział między dwie metody.

Metoda \texttt{CheckTx} jest tanim filtrem mempoola, niezmieniającym stanu. Wykonuje kolejno: dekodowanie transakcji, walidację składniową, kontrolę zgodności identyfikatora łańcucha, weryfikację podpisu oraz lekkie sprawdzenie semantyczne. To ostatnie celowo \emph{nie} egzekwuje dokładnej wartości licznika, a jedynie odrzuca transakcje, które nie mogą stać się poprawne (rejestracja istniejącej domeny, mutacja z licznikiem już przeterminowanym) -- ponieważ mempool może przechowywać transakcje w kolejności nieco innej niż docelowa, transakcja wyprzedzająca swoją poprzedniczkę nie powinna być przedwcześnie odrzucona. Dla transakcji rejestracji \texttt{CheckTx} wywołuje dodatkowo weryfikację własności domeny; ponieważ wiąże się ona z niedeterministycznym zapytaniem DNS, jest wykonywana \emph{poza} blokadą stanu i stanowi wyłącznie bramkę przyjęcia do mempoola.

Metoda \texttt{FinalizeBlock} jest autorytatywna i w pełni deterministyczna. Wykonuje każdą transakcję bloku po kolei, a wynik -- skrót stanu (korzeń drzewa) -- zwraca jako pole \texttt{AppHash} odpowiedzi. W ABCI 2.0 to właśnie tutaj, a nie w \texttt{Commit}, zwracany jest skrót stanu; sama metoda \texttt{Commit} jest w implementacji operacją pustą. Powtarza ona kontrole z \texttt{CheckTx}, lecz egzekwuje już \emph{dokładną} kolejną wartość licznika i faktycznie modyfikuje drzewo. Transakcja nieprawidłowa otrzymuje niezerowy kod wyniku i pozostawia stan nietknięty, lecz \emph{nie} przerywa przetwarzania bloku -- konsensus zgodził się już na włączenie tych bajtów, a transakcja odrzucona jest wynikiem przetwarzania, a nie awarią węzła. Determinizm tej ścieżki jest warunkiem koniecznym zgodności skrótów stanu między walidatorami; jego empiryczną weryfikację opisano w sekcji~\ref{sec:testy} oraz w rozdziale~\ref{ch:eksperymenty}.

Weryfikacja własności domeny powtarzana jest jeszcze w metodzie \texttt{ProcessProposal}, którą każdy walidator wykonuje na proponowanym bloku przed oddaniem głosu. Gdyby weryfikację umieścić wyłącznie w \texttt{CheckTx}, nieuczciwy proponent mógłby pominąć mempool i umieścić w bloku nieuwierzytelnioną rejestrację; ponowna weryfikacja na poziomie propozycji zamyka tę lukę, odrzucając cały blok zawierający rejestrację bez ważnego dowodu kontroli nad domeną. Podobnie jak \texttt{CheckTx}, \texttt{ProcessProposal} wykonuje niedeterministyczne zapytania DNS i z tego względu pozostaje bramką przedkonsensusową, poza deterministycznym \texttt{FinalizeBlock}.

\subsection{Zarządzanie stanem}
\label{subsec:zarzadzanie-stanem}

Stan aplikacji jest stosowany do drzewa Merkle \emph{bezpośrednio}, bez utrzymywania osobnej kopii roboczej. Decyzja ta wynika z użytego w prototypie magazynu w pamięci operacyjnej, który nie udostępnia tanich migawek; przy magazynie trwałym można by rozważyć kopię roboczą i atomowe zatwierdzanie. W obecnej konstrukcji skrót stanu jest po prostu odczytywany z drzewa na końcu \texttt{FinalizeBlock} i zwracany silnikowi konsensusu.

Konsekwencją braku trwałości jest utrata stanu przy zatrzymaniu węzła. Odzyskiwanie opiera się na mechanizmie wbudowanym w CometBFT: silnik przechowuje na dysku magazyn bloków, a po restarcie -- na podstawie wysokości i skrótu stanu zaraportowanych przez aplikację w odpowiedzi \texttt{Info} -- odtwarza brakujące bloki, ponownie wywołując \texttt{FinalizeBlock}. Ponieważ wykonanie jest deterministyczne, odtworzenie historii rekonstruuje dokładnie ten sam stan i ten sam skrót. Co istotne, ponieważ weryfikacja DNS jest bramką przedkonsensusową, a nie częścią \texttt{FinalizeBlock}, odtwarzanie nie powtarza zapytań DNS i przebiega poprawnie nawet wtedy, gdy rekordy TXT użyte przy pierwotnej rejestracji już nie istnieją. Skuteczność odzyskiwania -- powrót zrestartowanego walidatora do zgodnego stanu -- zweryfikowano empirycznie na sieci testowej (rozdział~\ref{ch:eksperymenty}).

Odwzorowanie domeny na stan realizuje pakiet \texttt{internal/state}: kluczem drzewa jest skrót SHA-256 nazwy domeny, a wartością liścia -- deterministyczna serializacja struktury \texttt{DomainState}. Magazyn, na którym osadzone jest drzewo, ukryto za wąskim interfejsem (operacje pobrania, zapisu, usunięcia i zamknięcia), co pozwala zastąpić implementację w pamięci magazynem trwałym bez zmiany logiki drzewa.

\subsection{Implementacja drzewa Merkle}
\label{subsec:implementacja-merkle}

Rzadkie drzewo Merkle (sekcja~\ref{sec:merkle}) zaimplementowano jako drzewo binarne o stałej głębokości 256, odpowiadającej długości klucza w bitach. Naiwna reprezentacja takiego drzewa byłaby niewykonalna -- liczba węzłów rzędu $2^{256}$ przekracza możliwości jakiegokolwiek magazynu. Wykorzystano zatem obserwację, że przygniatająca większość poddrzew jest pusta i ma z góry znane skróty. W fazie inicjalizacji wyznacza się tablicę \emph{skrótów domyślnych}: skrót pustego liścia to ciąg bajtów zerowych, a skrót pustego węzła na każdym wyższym poziomie powstaje przez połączenie dwóch identycznych skrótów domyślnych z poziomu niższego. Dzięki temu w magazynie przechowywane są wyłącznie węzły na ścieżkach kluczy faktycznie obecnych; aktualizacja jednego klucza dotyka co najwyżej 256 węzłów ścieżki od liścia do korzenia, a węzły, których skrót zapada się do wartości domyślnej, są z magazynu usuwane.

Szczególną uwagę poświęcono zwięzłości dowodów, od której zależy koszt po stronie strony ufającej. Dowód mógłby naiwnie zawierać 256 skrótów rodzeństwa -- po jednym na poziom -- jednak większość z nich to skróty domyślne, możliwe do odtworzenia przez weryfikującego na podstawie samego numeru poziomu. Zastosowano więc kompresję: struktura dowodu (\texttt{Proof}) niesie klucz, wartość (niepustą dla dowodu włączenia, pustą dla dowodu nieobecności), wektor wyłącznie \emph{niedomyślnych} skrótów rodzeństwa oraz 32-bajtową mapę bitową wskazującą, na których poziomach rodzeństwo jest niedomyślne. Weryfikujący odtwarza pozostałe skróty z tablicy domyślnych i przelicza ścieżkę do korzenia. Empiryczny rozmiar tak skompresowanego dowodu oraz jego zależność od liczby domen przebadano w rozdziale~\ref{ch:eksperymenty}; potwierdza on przewidywany wzrost rzędu $O(\log N)$. Identyczna postać dowodów włączenia i nieobecności sprawia, że ta sama procedura weryfikacji obsługuje oba przypadki, a dowód jest weryfikowalny wyłącznie względem korzenia, bez dostępu do reszty drzewa.

\section{Klient i CLI}
\label{sec:klient-cli}

Klient \texttt{dpki-cli} jest narzędziem wiersza poleceń zbudowanym na bibliotece Cobra. Przyjęto w nim nadrzędną zasadę: klient \emph{nie zawiera logiki biznesowej} -- każda decyzja walidacyjna należy do węzła, a rolą klienta jest jedynie zbudować transakcję, podpisać ją, wysłać i odczytać wynik bądź zweryfikować dowód. Dzięki temu reguły rejestru pozostają w jednym miejscu (aplikacja ABCI), a klient nie może stać się źródłem rozbieżnej interpretacji stanu.

Klient udostępnia komendy odpowiadające operacjom z sekcji~\ref{sec:protokoly} oraz operacjom pomocniczym:

\begin{itemize}
    \item \texttt{keys} -- generowanie i obsługa par kluczy Ed25519;
    \item \texttt{challenge} oraz \texttt{challenge-check} -- wyznaczenie wartości wyzwania DNS-01 do opublikowania oraz sprawdzenie jej obecności w DNS przed wysłaniem rejestracji;
    \item \texttt{register}, \texttt{rotate}, \texttt{revoke} -- budowa, podpisanie i rozgłoszenie odpowiednich transakcji;
    \item \texttt{lookup} -- odczyt stanu domeny wraz z dowodem i jego lokalną weryfikacją;
    \item \texttt{verify} -- weryfikacja w trybie wielowęzłowym oraz offline (sekcja~\ref{sec:api});
    \item \texttt{status} -- odczyt informacji o stanie węzła.
\end{itemize}

Komunikację z węzłem wydzielono do osobnego pakietu \texttt{client}, stanowiącego nakładkę na interfejs RPC silnika konsensusu. Hermetyzuje on rozgłaszanie transakcji oraz dwie ścieżki zapytań (o stan oraz o stan z dowodem), a także dekodowanie dowodu z odpowiedzi i pozyskanie korzenia zaufania z nagłówka kolejnego bloku. Logika weryfikacji dowodu jest współdzielona z aplikacją przez pakiet \texttt{internal/state}, co gwarantuje, że klient weryfikuje dowód dokładnie tą samą procedurą, którą węzeł go wytworzył.

\section{Konfiguracja sieci wielowęzłowej}
\label{sec:siec-wielowezlowa}

Aby umożliwić demonstrację właściwości rozproszonych systemu -- replikacji stanu i odporności na awarie -- zaimplementowano podkomendę \texttt{dpkid testnet init}, generującą kompletną konfigurację lokalnej sieci wielowęzłowej jednym wywołaniem. Rozwiązanie to świadomie zrealizowano w języku Go, a nie w postaci skryptu powłoki: generator korzysta z tych samych bibliotek konfiguracji, kluczy prywatnych walidatora i tożsamości sieciowej, których używa sam węzeł, dzięki czemu nie wymaga zewnętrznych narzędzi i eliminuje klasę błędów wynikających z ręcznej manipulacji plikami konfiguracyjnymi.

Generator tworzy katalogi domowe żądanej liczby walidatorów, a dla każdego z nich: unikatową parę kluczy konsensusu oraz klucz tożsamości sieciowej, wspólny dla całej sieci plik genezy (identyczny co do bajta -- najmniejsza różnica spowodowałaby odrzucenie węzła jako należącego do innego łańcucha) wymieniający wszystkich walidatorów oraz konfigurację z rozłącznymi portami. Węzeł $i$ nasłuchuje na porcie $26656 + 10i$ dla komunikacji P2P oraz $26657 + 10i$ dla RPC. Listę węzłów sąsiednich (\texttt{persistent\_peers}) zestawia się w pełną siatkę z rzeczywistych identyfikatorów tożsamości sieciowej; ponieważ na pojedynczej maszynie wszystkie węzły dzielą adres pętli zwrotnej, włączono opcje zezwalające na wielokrotne wystąpienie tego samego adresu. Tryb \texttt{-{}-docker} zmienia adresowanie węzłów sąsiednich na nazwy hostów \texttt{node0}, \texttt{node1} itd., umożliwiając uruchomienie sieci w kontenerach.

Zgodnie z arytmetyką konsensusu BFT (sekcja~\ref{subsec:bft-tendermint}) jako konfigurację referencyjną przyjęto cztery węzły -- najmniejszą liczbę dającą nietrywialną odporność na awarie ($n = 3f + 1$ dla $f = 1$). Uruchamianie i zatrzymywanie sieci ujęto w skryptach pomocniczych, a jej zachowanie w poszczególnych scenariuszach awarii przebadano w rozdziale~\ref{ch:eksperymenty}.

\section{Testy jednostkowe i integracyjne}
\label{sec:testy}

Implementację objęto wielowarstwowym zestawem testów, którego celem było nie tylko wykrywanie regresji, lecz także empiryczne potwierdzenie własności bezpieczeństwa, na których opiera się projekt. Wszystkie testy uruchamiane są z włączonym wykrywaniem wyścigów, co weryfikuje poprawność synchronizacji współbieżnego dostępu do stanu.

\begin{itemize}
    \item \textbf{testy jednostkowe} pokrywają poszczególne pakiety: prymitywy kryptograficzne, walidację i serializację modelu danych, operacje drzewa Merkle wraz z kodowaniem dowodów oraz weryfikację DNS-01;
    \item \textbf{testy integracyjne} weryfikują pełny przepływ przez aplikację ABCI -- od rejestracji, przez rotację i unieważnienie, po wygenerowanie dowodu i jego weryfikację offline -- a także sprawdzają determinizm wykonania, porównując skróty stanu uzyskane przy powtórzonym przetwarzaniu tej samej sekwencji transakcji;
    \item \textbf{testy bezpieczeństwa} odwzorowują scenariusze ataku z modelu adwersarza (sekcja~\ref{sec:zalozenia-adwersarz}) na konkretne przypadki testowe, stanowiąc empiryczną podstawę analizy w rozdziale~\ref{ch:analiza-bezpieczenstwa};
    \item \textbf{testy rozmyte} drzewa Merkle sprawdzają niezmienniki struktury dla losowo generowanych ciągów operacji.
\end{itemize}

Do testów wykorzystano dwa mechanizmy zastępcze, izolujące przedmiot testu od niedeterministycznych zależności. Pierwszym jest atrapa weryfikatora własności domeny, pozwalająca testować logikę rejestru bez rzeczywistych zapytań DNS; testy samego weryfikatora używają z kolei lokalnego serwera DNS odpowiadającego na zapytania w sposób kontrolowany. Drugim jest atrapa węzła wystawiająca interfejs RPC ponad rzeczywistą aplikacją, dzięki czemu testy klienta przebiegają na prawdziwych dowodach i skrótach stanu, lecz bez uruchamiania silnika konsensusu.

Pokrycie kodu testami, zmierzone narzędziem \texttt{go test}, przedstawia tabela~\ref{tab:pokrycie}. Najwyższe pokrycie osiągają pakiety o największym znaczeniu dla bezpieczeństwa -- weryfikator (100\%), prymitywy kryptograficzne, warstwa stanu oraz aplikacja ABCI. Niższa wartość dla pakietu \texttt{internal/types} wynika z wliczenia do niej kodu wygenerowanego automatycznie z Protocol Buffers, którego liczne metody pomocnicze nie są bezpośrednio wywoływane.

\begin{table}[htbp]
    \centering
    \caption{Pokrycie kodu testami w kluczowych pakietach}
    \label{tab:pokrycie}
    \begin{tabular}{l r}
        \hline
        \textbf{Pakiet}                & \textbf{Pokrycie} \\
        \hline
        \texttt{internal/verifier}     & 100,0\%           \\
        \texttt{internal/crypto}       & 92,3\%            \\
        \texttt{internal/state}        & 89,1\%            \\
        \texttt{internal/app}          & 86,7\%            \\
        \texttt{cmd/dpki-cli/commands} & 80,6\%            \\
        \hline
    \end{tabular}
\end{table}

Tak skonstruowany zestaw testów daje podstawę do dwóch kolejnych rozdziałów: testy bezpieczeństwa dostarczają materiału do analizy w rozdziale~\ref{ch:analiza-bezpieczenstwa}, a narzędzie pomiarowe \texttt{benchmark} oraz sieć testowa -- do oceny wydajnościowej w rozdziale~\ref{ch:eksperymenty}.
